\chapter{Related Work}
\label{sec:related_work}
% Your work needs to be grounded and compared to earlier work and the state-of-the-art. Start the section with announcing the research gap and also end with the research gap. Consider using hypotheses. 
Write about your related work here. Make clear to which key papers you will compare your eventual results. This can be done from the perspective of methods used, the task at hand and the addressed domain.

% The related work section is not a a background section. If you want to explain techniques, it is possible to have a background section after the Related Work section. Background should only be added when it has benefit for an informed audience. It is not common to use a background section. 

Remember: a related works section is not a ``laundry list'' of references, cited for the sake of citing papers: they are the support and motivation for your research gap and research question that you are going to address.

If you stick with Use \verb|\cite{}| for reference that is part of the sentence, and \verb|\citep{}| for references in parenthesis. For example, \cite{viola1997alignment} is awesome. Also, this is a citation~\citep{viola1997alignment}.